\begin{frame}
  \titlepage{}
\end{frame}

\section{\LaTeX{}: Was, warum, wie?}

\begin{frame}
    \frametitle{Was ist \LaTeX{}?}
    \begin{itemize}
        \item Eine logische Markup-Sprache zur Textauszeichnung
        \item Plattformunabhängig und Quelloffen
        \item Ideal für wissenschaftliche Dokumente
        \item Multilingual
    \end{itemize}
    \begin{description}[pdf\TeX]
        \item [\TeX] ist ein Textsatzsystem (1977--1986, D.\,E. Knuth)
        \item [\LaTeX] komfortable Benutzerschnittstelle (seit 1980, L.
        Lamport)
        \item [pdf\TeX] Ausgabesystem (seit 1990, H.\,T. Thành)
    \end{description}
    %TODO: Was ist Latex nicht
\end{frame}

\begin{frame}
    \frametitle{Warum \LaTeX?}
    \begin{itemize}
        \item Zeit- und Stressersparnis
        \item Plattformunabhängig
        \item Hohe typographische Qualität
        \item Gut dokumentiert (Internet, Bücher)
        \item Quasistandard bei wissenschaftlichen Veröffentlichungen
        \item Konzentration auf den Inhalt, nicht das Layout
    \end{itemize}
\end{frame}

\begin{frame}
    \frametitle{So funktioniert \LaTeX{}}
    \begin{itemize}
        \item Übersetzen von geschriebenem Quelltext in PDF-Dokument
        \item Programm dafür: pdflatex
        \item Für erweiterte Funktionen, z.B. Verweise mehrere Durchläufe nötig.
    \end{itemize}
    \pause\bigskip
    \begin{block}{Ablaufschema}\centering
        \begin{tikzpicture}[scale=.75,shorten <=2pt, shorten >=2pt]

            \node[rectangle,label=above:{Code}]
            	(code) 		at (0,0) {\includegraphics[height=30pt]{texicon.png}};
            \node[rectangle,draw,fill=white]
            	(latex) 	at (3,0) {\texttt{pdflatex}};
            \node[rectangle,label=above:{Dokument}]
            	(dokument) 	at (6,0) {\includegraphics[height=30pt]{pdficon.png}};
            \node
            	(aux) at (3,-2) {\footnotesize Hilfsdateien: \texttt{.aux}, \texttt{.toc},
            	\texttt{.log}, \ldots};
            \begin{scope}[thick]
                \path[->] (code)  edge (latex);
                \path[->] (latex) edge (dokument);
                \path[<->] (latex) edge (aux);
            \end{scope}
        \end{tikzpicture}
    \end{block}
\end{frame}
